\documentclass[aps,preprint]{revtex4}%
\usepackage{amssymb}
\usepackage{amsfonts}
\usepackage{amsmath}
\usepackage{graphicx}%
\setcounter{MaxMatrixCols}{30}
%TCIDATA{OutputFilter=latex2.dll}
%TCIDATA{Version=5.50.0.2960}
%TCIDATA{CSTFile=revtex4.cst}
%TCIDATA{Created=Saturday, July 31, 2010 09:28:33}
%TCIDATA{LastRevised=Friday, October 29, 2010 11:03:14}
%TCIDATA{<META NAME="GraphicsSave" CONTENT="32">}
%TCIDATA{<META NAME="SaveForMode" CONTENT="1">}
%TCIDATA{BibliographyScheme=Manual}
%TCIDATA{<META NAME="DocumentShell" CONTENT="Articles\SW\REVTeX 4">}
%BeginMSIPreambleData
\providecommand{\U}[1]{\protect\rule{.1in}{.1in}}
%EndMSIPreambleData
\newtheorem{theorem}{Theorem}
\newtheorem{acknowledgement}[theorem]{Acknowledgement}
\newtheorem{algorithm}[theorem]{Algorithm}
\newtheorem{axiom}[theorem]{Axiom}
\newtheorem{claim}[theorem]{Claim}
\newtheorem{conclusion}[theorem]{Conclusion}
\newtheorem{condition}[theorem]{Condition}
\newtheorem{conjecture}[theorem]{Conjecture}
\newtheorem{corollary}[theorem]{Corollary}
\newtheorem{criterion}[theorem]{Criterion}
\newtheorem{definition}[theorem]{Definition}
\newtheorem{example}[theorem]{Example}
\newtheorem{exercise}[theorem]{Exercise}
\newtheorem{lemma}[theorem]{Lemma}
\newtheorem{notation}[theorem]{Notation}
\newtheorem{problem}[theorem]{Problem}
\newtheorem{proposition}[theorem]{Proposition}
\newtheorem{remark}[theorem]{Remark}
\newtheorem{solution}[theorem]{Solution}
\newtheorem{summary}[theorem]{Summary}
\newenvironment{proof}[1][Proof]{\noindent\textbf{#1.} }{\ \rule{0.5em}{0.5em}}
\begin{document}
\preprint{ }
\title{Berezin Transform}
\author{Brian Weiner}
\affiliation{Penn State Univ}

\begin{abstract}


\end{abstract}
\date{08/29/2010}
\startpage{1}
\maketitle
\tableofcontents


\section{Operator Kernels}

Using
\begin{equation}
I=\int\left\vert \mathbf{z}\right\rangle \left\langle \mathbf{z}\right\vert
d\mu\left(  \mathbf{z}\right)  =\int\frac{\left\vert \mathbf{z}\right\rangle
\left\langle \mathbf{z}\right\vert }{\left\langle \mathbf{z}\left\vert
\mathbf{z}\right.  \right\rangle }d\tilde{\mu}\left(  \mathbf{z}\right)  ,
\end{equation}
where
\begin{equation}
d\tilde{\mu}\left(  \mathbf{z}\right)  =\left\langle \mathbf{z}\left\vert
\mathbf{z}\right.  \right\rangle d\mu\left(  \mathbf{z}\right)
\end{equation}
and
\begin{equation}
A=IAI
\end{equation}
one can construct a biholomorphic representations of an operator $A$ wrt to
the measures $\mu$ and $\tilde{\mu}$ by
\begin{align}
A &  =\int\left\vert \mathbf{z}_{1}\right\rangle \left\langle \mathbf{z}%
_{1}\right\vert A\left\vert \mathbf{z}_{2}\right\rangle \left\langle
\mathbf{z}_{2}\right\vert d\mu\left(  \mathbf{z}_{1}\right)  d\mu\left(
\mathbf{z}_{2}\right)  \nonumber\\
&  =\int\frac{\left\vert \mathbf{z}_{1}\right\rangle }{\sqrt{\left\langle
\mathbf{z}_{1}\left\vert \mathbf{z}_{1}\right.  \right\rangle }}%
\frac{\left\langle \mathbf{z}_{1}\left\vert A\mathbf{z}_{2}\right.
\right\rangle }{\sqrt{\left\langle \mathbf{z}_{1}\left\vert \mathbf{z}%
_{1}\right.  \right\rangle \left\langle \mathbf{z}_{2}\left\vert
\mathbf{z}_{2}\right.  \right\rangle }}\frac{\left\langle \mathbf{z}%
_{2}\right\vert }{\sqrt{\left\langle \mathbf{z}_{2}\left\vert \mathbf{z}%
_{2}\right.  \right\rangle }}d\tilde{\mu}\left(  \mathbf{z}_{1}\right)
d\tilde{\mu}\left(  \mathbf{z}_{2}\right)  \nonumber\\
&  =\int A\left(  \mathbf{\bar{z}}_{1},\mathbf{z}_{2}\right)  \left\vert
\mathbf{z}_{1}\right\rangle \left\langle \mathbf{z}_{2}\right\vert d\mu\left(
\mathbf{z}_{1}\right)  d\mu\left(  \mathbf{z}_{2}\right)  \nonumber\\
&  =\int\tilde{A}\left(  \mathbf{\bar{z}}_{1},\mathbf{z}_{2}\right)
\frac{\left\vert \mathbf{z}_{1}\right\rangle \left\langle \mathbf{z}%
_{2}\right\vert }{\sqrt{\left\langle \mathbf{z}_{1}\left\vert \mathbf{z}%
_{1}\right.  \right\rangle \left\langle \mathbf{z}_{2}\left\vert
\mathbf{z}_{2}\right.  \right\rangle }}d\tilde{\mu}\left(  \mathbf{z}%
_{1}\right)  d\tilde{\mu}\left(  \mathbf{z}_{2}\right)  ,
\end{align}
where
\begin{subequations}
\label{Q-Symbols}%
\begin{align}
A\left(  \mathbf{\bar{z}}_{1},\mathbf{z}_{2}\right)   &  =\left\langle
\mathbf{z}_{1}\left\vert A\mathbf{z}_{2}\right.  \right\rangle \nonumber\\
\tilde{A}\left(  \mathbf{\bar{z}}_{1},\mathbf{z}_{2}\right)   &
=\frac{\left\langle \mathbf{z}_{1}\left\vert A\mathbf{z}_{2}\right.
\right\rangle }{\sqrt{\left\langle \mathbf{z}_{1}\left\vert \mathbf{z}%
_{1}\right.  \right\rangle \left\langle \mathbf{z}_{2}\left\vert
\mathbf{z}_{2}\right.  \right\rangle }}.
\end{align}
These representations have the following action on a holomorphic wavefunction
\end{subequations}
\begin{align}
A\left\vert g\right\rangle  &  =\left\vert Ag\right\rangle =\nonumber\\
\int\left\vert \mathbf{z}_{1}\right\rangle \left\langle \mathbf{z}%
_{1}\right\vert A\left\vert \mathbf{z}_{2}\right\rangle \left\langle
\mathbf{z}_{2}\right\vert \left\vert g\right\rangle d\mu\left(  \mathbf{z}%
_{1}\right)  d\mu\left(  \mathbf{z}_{2}\right)   &  =\int A\left(
\mathbf{\bar{z}}_{1},\mathbf{z}_{2}\right)  \left\vert \mathbf{z}%
_{1}\right\rangle \left\langle \mathbf{z}_{2}\right\vert \left\vert
g\right\rangle d\mu\left(  \mathbf{z}_{1}\right)  d\mu\left(  \mathbf{z}%
_{2}\right)  \nonumber\\
&  =\int\left(  A\left(  \mathbf{\bar{z}}_{1},\mathbf{z}_{2}\right)
\left\langle \mathbf{z}_{2}\right\vert \left\vert g\right\rangle d\mu\left(
\mathbf{z}_{2}\right)  \right)  \left\vert \mathbf{z}_{1}\right\rangle
d\mu\left(  \mathbf{z}_{1}\right)  \nonumber\\
&  =\int\left(  A\left(  \mathbf{\bar{z}}_{1},\mathbf{z}_{2}\right)  g\left(
\mathbf{\bar{z}}_{2}\right)  d\mu\left(  \mathbf{z}_{2}\right)  \right)
\left\vert \mathbf{z}_{1}\right\rangle d\mu\left(  \mathbf{z}_{1}\right)
\nonumber\\
\left[  Ag\right]  \left(  \mathbf{\bar{z}}_{1}\right)   &  =\int\left(
A\left(  \mathbf{\bar{z}}_{1},\mathbf{z}_{2}\right)  g\left(  \mathbf{\bar{z}%
}_{2}\right)  d\mu\left(  \mathbf{z}_{2}\right)  \right)
\end{align}
and%
\begin{align}
A\left\vert g\right\rangle = &  \int\frac{\left\vert \mathbf{z}_{1}%
\right\rangle }{\sqrt{\left\langle \mathbf{z}_{1}\left\vert \mathbf{z}%
_{1}\right.  \right\rangle }}\frac{\left\langle \mathbf{z}_{1}\left\vert
A\mathbf{z}_{2}\right.  \right\rangle }{\sqrt{\left\langle \mathbf{z}%
_{1}\left\vert \mathbf{z}_{1}\right.  \right\rangle \left\langle
\mathbf{z}_{2}\left\vert \mathbf{z}_{2}\right.  \right\rangle }}%
\frac{\left\langle \mathbf{z}_{2}\right\vert }{\sqrt{\left\langle
\mathbf{z}_{2}\left\vert \mathbf{z}_{2}\right.  \right\rangle }}\left\vert
g\right\rangle d\tilde{\mu}\left(  \mathbf{z}_{1}\right)  d\tilde{\mu}\left(
\mathbf{z}_{2}\right)  \nonumber\\
&  =\int\left(  \tilde{A}\left(  \mathbf{\bar{z}}_{1},\mathbf{z}_{2}\right)
\tilde{g}\left(  \mathbf{\bar{z}}_{2}\right)  d\tilde{\mu}\left(
\mathbf{z}_{2}\right)  \right)  \frac{\left\vert \mathbf{z}_{1}\right\rangle
}{\sqrt{\left\langle \mathbf{z}_{1}\left\vert \mathbf{z}_{1}\right.
\right\rangle }}d\tilde{\mu}\left(  \mathbf{z}_{1}\right)  \nonumber\\
\widetilde{\left[  Ag\right]  }\left(  \mathbf{\bar{z}}_{1}\right)   &
=\int\tilde{A}\left(  \mathbf{\bar{z}}_{1},\mathbf{z}_{2}\right)  \tilde
{g}\left(  \mathbf{\bar{z}}_{2}\right)  d\tilde{\mu}\left(  \mathbf{z}%
_{2}\right)  .
\end{align}


The ket-bra expansion of an operator $A$ in the reference basis $\left\{
\left.  \left\vert \varphi_{j_{1}}\varphi_{j_{2}}\right\rangle \left\langle
\varphi_{k_{1}}\varphi_{k_{2}}\right\vert \right\vert 1\leq j_{1}<j_{2}%
,k_{1}<k_{2}\leq r\right\}  $ is
\begin{align}
A  &  =\sum_{1\leq j_{1}<j_{2},k_{1}<k_{2}\leq r}\left\vert \varphi_{j_{1}%
}\varphi_{j_{2}}\right\rangle \left\langle \varphi_{j_{1}}\varphi_{j_{2}%
}\right\vert A\left\vert \varphi_{k_{1}}\varphi_{k_{2}}\right\rangle
\left\langle \varphi_{k_{1}}\varphi_{k_{2}}\right\vert \nonumber\\
&  =\sum_{1\leq j_{1}<j_{2},k_{1}<k_{2}\leq r}A_{j_{1}j_{2}k_{1}k_{2}%
}\left\vert \varphi_{j_{1}}\varphi_{j_{2}}\right\rangle \left\langle
\varphi_{k_{1}}\varphi_{k_{2}}\right\vert ,
\end{align}
where
\begin{equation}
A_{j_{1}j_{2}k_{1}k_{2}}=\operatorname*{Tr}\left\{  \left\vert \varphi_{k_{1}%
}\varphi_{k_{2}}\right\rangle \left\langle \varphi_{j_{1}}\varphi_{j_{2}%
}\right\vert A\right\}  =\operatorname*{Tr}\left\{  \left\vert \varphi_{j_{1}%
}\varphi_{j_{2}}\right\rangle \left\langle \varphi_{k_{1}}\varphi_{k_{2}%
}\right\vert ^{\dagger}A\right\}  .
\end{equation}
The biholomorphic integral kernel $A$ can thus be expanded as
\begin{align}
A\left(  \mathbf{\bar{z}}_{1},\mathbf{z}_{2}\right)   &  =\sum_{1\leq
j_{1}<j_{2},k_{1}<k_{2}\leq r}A_{j_{1}j_{2}k_{1}k_{2}}\left\langle
\mathbf{z}_{1}\left\vert \varphi_{j_{1}}\varphi_{j_{2}}\right.  \right\rangle
\left\langle \varphi_{k_{1}}\varphi_{k_{2}}\left\vert \mathbf{z}_{2}\right.
\right\rangle \nonumber\\
&  =\sum_{1\leq j_{1}<j_{2},k_{1}<k_{2}\leq r}A_{j_{1}j_{2}k_{1}k_{2}%
}\mathcal{T}_{k_{1}k_{2}j_{1}j_{2}}\left(  \mathbf{\bar{z}}_{1},\mathbf{z}%
_{2}\right) \nonumber\\
&  =\sum_{1\leq j_{1}<j_{2},k_{1}<k_{2}\leq r}A_{j_{1}j_{2}k_{1}k_{2}%
}\mathcal{\tilde{T}}_{k_{1}k_{2}j_{1}j_{2}}\left(  \mathbf{\bar{z}}%
_{1},\mathbf{z}_{2}\right)  \sqrt{\left\langle \mathbf{z}_{1}\left\vert
\mathbf{z}_{1}\right.  \right\rangle \left\langle \mathbf{z}_{2}\left\vert
\mathbf{z}_{2}\right.  \right\rangle },
\end{align}
where the biholomorphic integral kernels of the bases are related by
\begin{equation}
\mathcal{T}_{j_{1}j_{2}k_{1}k_{2}}\left(  \mathbf{\bar{z}}_{1},\mathbf{z}%
_{2}\right)  =\mathcal{\tilde{T}}_{j_{1}j_{2}k_{1}k_{2}}\left(  \mathbf{\bar
{z}}_{1},\mathbf{z}_{2}\right)  \sqrt{\left\langle \mathbf{z}_{1}\left\vert
\mathbf{z}_{1}\right.  \right\rangle \left\langle \mathbf{z}_{2}\left\vert
\mathbf{z}_{2}\right.  \right\rangle }%
\end{equation}
and
\begin{subequations}
\label{Q-Sym}%
\begin{align}
\mathcal{\tilde{T}}_{j_{1}j_{2}k_{1}k_{2}}\left(  \mathbf{\bar{z}}%
_{1},\mathbf{z}_{2}\right)   &  =\frac{\operatorname*{Tr}\left\{  \left\vert
\varphi_{k_{1}}\varphi_{k_{2}}\right\rangle \left\langle \varphi_{j_{1}%
}\varphi_{j_{2}}\right\vert \left\vert \mathbf{z}_{2}\right\rangle
\left\langle \mathbf{z}_{1}\right\vert \right\}  }{\sqrt{\left\langle
\mathbf{z}_{1}\left\vert \mathbf{z}_{1}\right.  \right\rangle \left\langle
\mathbf{z}_{2}\left\vert \mathbf{z}_{2}\right.  \right\rangle }}\nonumber\\
\mathcal{T}_{j_{1}j_{2}k_{1}k_{2}}\left(  \mathbf{\bar{z}}_{1},\mathbf{z}%
_{2}\right)   &  =\operatorname*{Tr}\left\{  \left\vert \varphi_{k_{1}}%
\varphi_{k_{2}}\right\rangle \left\langle \varphi_{j_{1}}\varphi_{j_{2}%
}\right\vert \left\vert \mathbf{z}_{2}\right\rangle \left\langle
\mathbf{z}_{1}\right\vert \right\}  .
\end{align}


The kernels $\left\{  \left.  \mathcal{T}_{j_{1}j_{2}k_{1}k_{2}}\left(
\mathbf{\bar{z}}_{1},\mathbf{z}_{2}\right)  \right\vert 1\leq j_{1}%
<j_{2},k_{1}<k_{2}\leq r\right\}  $ form a biholomorphic conb for
$\mathfrak{F}_{op}=\mathcal{R}_{hol}\left(  \mathcal{B}_{2}\left(
\mathcal{H}^{2}\right)  \right)  \subset L_{2}\left(  \mathbb{C}^{4\left(
r-2\right)  },\mu\right)  $ as%
\end{subequations}
\begin{align}
&  \int\overline{\mathcal{T}_{j_{1}j_{2}k_{1}k_{2}}\left(  \mathbf{\bar{z}%
}_{1},\mathbf{z}_{2}\right)  }\mathcal{T}_{p_{1}p_{2}q_{1}q_{2}}\left(
\mathbf{\bar{z}}_{1},\mathbf{z}_{2}\right)  d\mu\left(  \mathbf{z}_{1}\right)
d\mu\left(  \mathbf{z}_{2}\right) \nonumber\\
&  =\int\overline{\left\langle \mathbf{z}_{1}\left\vert \varphi_{j_{1}}%
\varphi_{j_{2}}\right.  \right\rangle \left\langle \varphi_{k_{1}}%
\varphi_{k_{2}}\left\vert \mathbf{z}_{2}\right.  \right\rangle }\left\langle
\mathbf{z}_{1}\left\vert \varphi_{p_{1}}\varphi_{p_{2}}\right.  \right\rangle
\left\langle \varphi_{q_{1}}\varphi_{q_{2}}\left\vert \mathbf{z}_{2}\right.
\right\rangle d\mu\left(  \mathbf{z}_{1}\right)  d\mu\left(  \mathbf{z}%
_{2}\right) \nonumber\\
&  =\int\left\langle \varphi_{q_{1}}\varphi_{q_{2}}\left\vert \mathbf{z}%
_{2}\right.  \right\rangle \left\langle \mathbf{z}_{2}\left\vert
\varphi_{k_{1}}\varphi_{k_{2}}\right.  \right\rangle \left\langle
\varphi_{j_{1}}\varphi_{j_{2}}\left\vert \mathbf{z}_{1}\right.  \right\rangle
\left\langle \mathbf{z}_{1}\left\vert \varphi_{p_{1}}\varphi_{p_{2}}\right.
\right\rangle d\mu\left(  \mathbf{z}_{1}\right)  d\mu\left(  \mathbf{z}%
_{2}\right) \nonumber\\
&  =\left\langle \varphi_{q_{1}}\varphi_{q_{2}}\left\vert \varphi_{k_{1}%
}\varphi_{k_{2}}\right.  \right\rangle \left\langle \varphi_{j_{1}}%
\varphi_{j_{2}}\left\vert \varphi_{p_{1}}\varphi_{p_{2}}\right.  \right\rangle
=\delta_{\left(  q_{1}q_{2}\right)  \left(  k_{1}k_{2}\right)  }%
\delta_{\left(  j_{1}j_{2}\right)  \left(  p_{1}p_{2}\right)  }.
\end{align}


\subsection{Q-Symbols}

The Q-symbols, $\left\{  Q_{j_{1}j_{2}k_{1}k_{2}}\right\}  $ and $\left\{
\tilde{Q}_{j_{1}j_{2}k_{1}k_{2}}\right\}  ,$ wrt the measures $\mu$ and
$\tilde{\mu}$ of the ket-bra basis elements are given by
\begin{equation}
Q_{j_{1}j_{2}k_{1}k_{2}}\left(  \mathbf{z}\right)  =\operatorname*{Tr}\left\{
\left\vert \varphi_{k_{1}}\varphi_{k_{2}}\right\rangle \left\langle
\varphi_{j_{1}}\varphi_{j_{2}}\right\vert \left\vert \mathbf{z}\right\rangle
\left\langle \mathbf{z}\right\vert \right\}  =\tilde{Q}_{j_{1}j_{2}k_{1}k_{2}%
}\left(  \mathbf{z}\right)  \left\langle \mathbf{z}\left\vert \mathbf{z}%
\right.  \right\rangle
\end{equation}
and thus the Q-symbol, $Q_{A}\left(  \mathbf{z}\right)  $ and $\tilde{Q}%
_{A}\left(  \mathbf{z}\right)  ,$ of the operator $A$ by
\begin{equation}
Q_{A}\left(  \mathbf{z}\right)  =\tilde{Q}_{A}\left(  \mathbf{z}\right)
\left\langle \mathbf{z}\left\vert \mathbf{z}\right.  \right\rangle
=\sum_{1\leq j_{1}<j_{2},k_{1}<k_{2}\leq r}A_{j_{1}j_{2}k_{1}k_{2}}\tilde
{Q}_{j_{1}j_{2}k_{1}k_{2}}\left(  \mathbf{z}\right)  \left\langle
\mathbf{z}\left\vert \mathbf{z}\right.  \right\rangle .
\end{equation}


The set $\left\{  \left.  Q_{j_{1}j_{2}k_{1}k_{2}}\left(  \mathbf{z}\right)
\right\vert 1\leq j_{1}<j_{2}\leq r,1\leq k_{1}<k_{2}\leq r\right\}  $ is
linearly independent and forms a basis for the space, $\mathfrak{F}\subset
L_{2}\left(  \mathbb{C}^{2\left(  r-2\right)  },\mu\right)  ,$ of Q-symbols,
whose dimension is $\binom{r}{2}^{2}$. However, in general these functions are
not holomorphic, biholomorphic, normalized nor orthogonal and their metric
matrix is given by
\begin{align}
&  \Delta_{l_{1}l_{2}m_{1}m_{2}j_{1}j_{2}k_{1}k_{2}}=\left(  Q_{l_{1}%
l_{2}m_{1}m_{2}}\left\vert Q_{j_{1}j_{2}k_{1}k_{2}}\right.  \right)
_{\mathfrak{F}}=\nonumber\\
&  =\int\overline{\operatorname*{Tr}\left\{  \left\vert \mathbf{z}%
\right\rangle \left\langle \mathbf{z}\right\vert \left\vert \varphi_{l_{1}%
}\varphi_{l_{2}}\right\rangle \left\langle \varphi_{m_{1}}\varphi_{m_{2}%
}\right\vert \right\}  }\operatorname*{Tr}\left\{  \left\vert \mathbf{z}%
\right\rangle \left\langle \mathbf{z}\right\vert \left\vert \varphi_{j_{1}%
}\varphi_{j_{2}}\right\rangle \left\langle \varphi_{k_{1}}\varphi_{k_{2}%
}\right\vert \right\}  d\mu\left(  \mathbf{z}\right) \nonumber\\
&  =\int_{\mathbb{C}^{2\left(  r-2\right)  }}\overline{\gamma_{2}\left(
\mathbf{z}\right)  }_{m_{1}m_{2}l_{1}l_{2}}\gamma_{2}\left(  \mathbf{z}%
\right)  _{j_{1}j_{2}k_{1}k_{2}}d\mu\left(  \mathbf{z}\right)  .
\end{align}


\subsection{P-Symbols}

One can form a (in general non unique) biorthogonal basis $\left\{  \left.
P_{j_{1}j_{2}k_{1}k_{2}}\left(  \mathbf{z}\right)  \right\vert 1\leq
j_{1}<j_{2},k_{1}<k_{2}\leq r\right\}  $ to $\left\{  \left.  Q_{j_{1}%
j_{2}k_{1}k_{2}}\left(  \mathbf{z}\right)  \right\vert 1\leq j_{1}<j_{2}%
,k_{1}<k_{2}\leq r\right\}  $ that also spans $\mathfrak{F},$ i.e.
\begin{align}
\left(  P_{j_{1}j_{2}k_{1}k_{2}}\left\vert Q_{l_{1}l_{2}m_{1}m_{2}}\right.
\right)  _{\mathfrak{F}}  &  =\int\overline{P_{j_{1}j_{2}k_{1}k_{2}}\left(
\mathbf{z}\right)  }Q_{j_{1}j_{2}k_{1}k_{2}}\left(  \mathbf{z}\right)
d\mu\left(  \mathbf{z}\right) \nonumber\\
&  =\delta_{\left(  l_{1}l_{2}\right)  \left(  j_{1}j_{2}\right)  }%
\delta_{\left(  m_{1}m_{2}\right)  \left(  k_{1}k_{2}\right)  }.
\end{align}
The $\left\{  P_{j_{1}j_{2}k_{1}k_{2}}\left(  \mathbf{z}\right)  \right\}  $
basis can be expanded in terms of the $\left\{  Q_{j_{1}j_{2}k_{1}k_{2}%
}\left(  \mathbf{z}\right)  \right\}  $ basis as
\begin{align}
P_{j_{1}j_{2}k_{1}k_{2}}\left(  \mathbf{z}\right)   &  =\sum_{1\leq
p_{1}<p_{2},q_{1}<q_{2}\leq r}Q_{p_{1}p_{2}q_{1}q_{2}}\left(  \mathbf{z}%
\right)  \Delta_{p_{1}p_{2}q_{1}q_{2}j_{1}j_{2}k_{1}k_{2}}^{-1}\nonumber\\
&  =\sum_{1\leq p_{1}<p_{2},q_{1}<q_{2}\leq r}\gamma_{2}\left(  \mathbf{z}%
\right)  _{p_{1}p_{2}q_{1}q_{2}}\Delta_{p_{1}p_{2}q_{1}q_{2}j_{1}j_{2}%
k_{1}k_{2}}^{-1},
\end{align}
i.e.
\begin{equation}
\left\vert P\right)  =\left\vert Q\right)  \Delta^{-1}.
\end{equation}
This leads to
\begin{equation}
\left(  \left.  P\right\vert P\right)  =\Delta^{-1}\left(  \left.
Q\right\vert Q\right)  \Delta^{-1}=\Delta^{-1},
\end{equation}%
\begin{align}
\left\vert Q\right)   &  =\left\vert P\right)  \Delta,\nonumber\\
\left\vert P\right)   &  =\left\vert Q\right)  \Delta^{-1},
\end{align}
and%
\begin{equation}
I_{\mathfrak{F}}=\left\vert Q\right)  \Delta^{-1}\left(  Q\right\vert
=\left\vert P\right)  \Delta\left(  P\right\vert =\left\vert P\right)  \left(
Q\right\vert =\left\vert Q\right)  \left(  P\right\vert .
\end{equation}


One can also observe that
\begin{equation}
I_{\mathfrak{F}}=\int\left\vert \left\vert \mathbf{z}\right\rangle
\left\langle \mathbf{z}\right\vert \right)  \left(  \left\vert \mathbf{z}%
\right\rangle \left\langle \mathbf{z}\right\vert \right\vert d\mu\left(
\mathbf{z}\right)  ,
\end{equation}
which comes directly from
\begin{align}
&  \Delta_{l_{1}l_{2}m_{1}m_{2}j_{1}j_{2}k_{1}k_{2}}=\left(  Q_{l_{1}%
l_{2}m_{1}m_{2}}\left\vert Q_{j_{1}j_{2}k_{1}k_{2}}\right.  \right)
=\nonumber\\
&  =\int_{\mathbb{C}^{2\left(  r-2\right)  }}\overline{\operatorname*{Tr}%
\left\{  \left\vert \mathbf{z}\right\rangle \left\langle \mathbf{z}\right\vert
\left\vert \varphi_{l_{1}}\varphi_{l_{2}}\right\rangle \left\langle
\varphi_{m_{1}}\varphi_{m_{2}}\right\vert \right\}  }\operatorname*{Tr}%
\left\{  \left\vert \mathbf{z}\right\rangle \left\langle \mathbf{z}\right\vert
\left\vert \varphi_{j_{1}}\varphi_{j_{2}}\right\rangle \left\langle
\varphi_{k_{1}}\varphi_{k_{2}}\right\vert \right\}  d\mu\left(  \mathbf{z}%
\right) \nonumber\\
&  =\int_{\mathbb{C}^{2\left(  r-2\right)  }}\left(  \left.  Q_{l_{1}%
l_{2}m_{1}m_{2}}\right\vert \left\vert \mathbf{z}\right\rangle \left\langle
\mathbf{z}\right\vert \right)  \left(  \left\vert \mathbf{z}\right\rangle
\left\langle \mathbf{z}\right\vert \left\vert Q_{j_{1}j_{2}k_{1}k_{2}}\right.
\right)  d\mu\left(  \mathbf{z}\right)  .
\end{align}
This orthogonality condition can be reexpressed as
\begin{align}
\delta_{\left(  l_{1}l_{2}\right)  \left(  j_{1}j_{2}\right)  }\delta_{\left(
m_{1}m_{2}\right)  \left(  k_{1}k_{2}\right)  }  &  =\int P_{j_{1}j_{2}%
k_{1}k_{2}}\left(  \mathbf{z}\right)  Q_{l_{1}l_{2}m_{1}m_{2}}\left(
\mathbf{z}\right)  d\mu\left(  \mathbf{z}\right) \nonumber\\
&  =\int P_{j_{1}j_{2}k_{1}k_{2}}\left(  \mathbf{z}\right)  \overline
{\operatorname*{Tr}\left\{  \left\vert \mathbf{z}\right\rangle \left\langle
\mathbf{z}\right\vert \left\vert \varphi_{l_{1}}\varphi_{l_{2}}\right\rangle
\left\langle \varphi_{m_{1}}\varphi_{m_{2}}\right\vert \right\}  }d\mu\left(
\mathbf{z}\right) \nonumber\\
&  =\operatorname*{Tr}\left\{  \left(  \int\left\vert \mathbf{z}\right\rangle
\left\langle \mathbf{z}\right\vert P_{j_{1}j_{2}k_{1}k_{2}}\left(
\mathbf{z}\right)  d\mu\left(  \mathbf{z}\right)  \right)  \left\vert
\varphi_{m_{1}}\varphi_{m_{2}}\right\rangle \left\langle \varphi_{l_{1}%
}\varphi_{l_{2}}\right\vert \right\}  ,
\end{align}
which shows that the integral can be interpreted as the following operator
\begin{equation}
\left\vert \varphi_{j_{1}}\varphi_{j_{2}}\right\rangle \left\langle
\varphi_{k_{1}}\varphi_{k_{2}}\right\vert \equiv\int P_{j_{1}j_{2}k_{1}k_{2}%
}\left(  \mathbf{z}\right)  \left\vert \mathbf{z}\right\rangle \left\langle
\mathbf{z}\right\vert d\mu\left(  \mathbf{z}\right)
\end{equation}
and $P_{j_{1}j_{2}k_{1}k_{2}}\left(  \mathbf{z}\right)  $ as its P-symbol ,
i.e
\begin{equation}
P_{j_{1}j_{2}k_{1}k_{2}}\left(  \mathbf{z}\right)  =P_{j_{1}j_{2}k_{1}k_{2}%
}\left(  \mathbf{z}\right)  ,
\end{equation}
note that P-symbols are not unique.

The P-symbol, $P_{A}\left(  \mathbf{z}\right)  ,$ of any operator $A$ can then
be expressed as%
\begin{equation}
P_{A}\left(  \mathbf{z}\right)  =\sum_{1\leq j_{1}<j_{2},k_{1}<k_{2}\leq
r}A_{j_{1}j_{2}k_{1}k_{2}}P_{j_{1}j_{2}k_{1}k_{2}}\left(  \mathbf{z}\right)  ,
\end{equation}
in particular
\begin{align}
P_{\mathcal{H}^{2}} &  =\sum_{1\leq j_{1}<j_{2}\leq r}\left\vert
\varphi_{j_{1}}\varphi_{j_{2}}\right\rangle \left\langle \varphi_{j_{1}%
}\varphi_{j_{2}}\right\vert =\sum_{1\leq j_{1}<j_{2}\leq r}\int P_{j_{1}%
j_{2}j_{1}j_{2}}\left(  \mathbf{z}\right)  \left\vert \mathbf{z}\right\rangle
\left\langle \mathbf{z}\right\vert d\mu\left(  \mathbf{z}\right)  \nonumber\\
&  =\int\sum_{1\leq j_{1}<j_{2}\leq r}P_{j_{1}j_{2}j_{1}j_{2}}\left(
\mathbf{z}\right)  \left\vert \mathbf{z}\right\rangle \left\langle
\mathbf{z}\right\vert d\mu\left(  \mathbf{z}\right)  =\int P_{P_{\mathcal{H}%
^{2}}}\left(  \mathbf{z}\right)  \left\vert \mathbf{z}\right\rangle
\left\langle \mathbf{z}\right\vert d\mu\left(  \mathbf{z}\right)  ,
\end{align}
thus%
\begin{equation}
P_{P_{\mathcal{H}^{2}}}\left(  \mathbf{z}\right)  \equiv P_{\mathcal{H}^{2}%
}\left(  \mathbf{z}\right)  =\sum_{1\leq j_{1}<j_{2}\leq r}P_{j_{1}j_{2}%
j_{1}j_{2}}\left(  \mathbf{z}\right)  =1.
\end{equation}
This gives%
\begin{equation}
\left\langle \mathbf{z}^{\prime\prime}\left\vert \mathbf{z}^{\prime}\right.
\right\rangle =\int\left\langle \mathbf{z}^{\prime\prime}\left\vert
\mathbf{z}\right.  \right\rangle \left\langle \mathbf{z}\left\vert
\mathbf{z}^{\prime}\right.  \right\rangle d\mu\left(  \mathbf{z}\right)
\end{equation}
and
\begin{equation}
Q_{P_{\mathcal{H}^{2}}}\left(  \mathbf{z}^{\prime}\right)  =\int\left\vert
\left\langle \mathbf{z}^{\prime}\left\vert \mathbf{z}\right.  \right\rangle
\right\vert ^{2}d\mu\left(  \mathbf{z}\right)  .
\end{equation}


\subsection{Reproducing Kernel and Berezin Transform}

The reproducing kernel of $\mathcal{B}_{2}\left(  \mathfrak{F}\right)  $ is
\begin{align}
M\left(  \mathbf{z,z}^{\prime}\right)   &  =\sum_{1\leq j_{1}<j_{2}%
,k_{1}<k_{2}\leq r}\left(  \left\vert \mathbf{z}\right\rangle \left\langle
\mathbf{z}\right\vert \left\vert P_{j_{1}j_{2}k_{1}k_{2}}\right.  \right)
\left(  \left.  Q_{j_{1}j_{2}k_{1}k_{2}}\right\vert \left\vert \mathbf{z}%
^{\prime}\right\rangle \left\langle \mathbf{z}^{\prime}\right\vert \right)
\nonumber\\
&  =\sum_{\substack{1\leq j_{1}<j_{2},k_{1}<k_{2}\leq r\\1\leq p_{1}%
<p_{2},q_{1}<q_{2}\leq r}}\left(  \left\vert \mathbf{z}\right\rangle
\left\langle \mathbf{z}\right\vert \left\vert Q_{j_{1}j_{2}k_{1}k_{2}}\right.
\right)  \Delta_{j_{1}j_{2}k_{1}k_{2}p_{1}p_{2}q_{1}q_{2}}^{-1}\left(  \left.
Q_{p_{1}p_{2}q_{1}q_{2}}\right\vert \left\vert \mathbf{z}^{\prime
}\right\rangle \left\langle \mathbf{z}^{\prime}\right\vert \right) \nonumber\\
&  =\sum_{\substack{1\leq j_{1}<j_{2},k_{1}<k_{2}\leq r\\1\leq p_{1}%
<p_{2},q_{1}<q_{2}\leq r}}\left(  \left\vert \mathbf{z}\right\rangle
\left\langle \mathbf{z}\right\vert \left\vert P_{j_{1}j_{2}k_{1}k_{2}}\right.
\right)  \Delta_{j_{1}j_{2}k_{1}k_{2}p_{1}p_{2}q_{1}q_{2}}\left(  \left.
P_{p_{1}p_{2}q_{1}q_{2}}\right\vert \left\vert \mathbf{z}^{\prime
}\right\rangle \left\langle \mathbf{z}^{\prime}\right\vert \right)  .
\end{align}


The Kernel of the Berezin Transformation that transform P-symbols to
Q-symbols
\begin{equation}
\mathfrak{Be}:\mathfrak{F}\rightarrow\mathfrak{F}%
\end{equation}
can be obtained by using the identity
\begin{equation}
\left\vert Q_{j_{1}j_{2}k_{1}k_{2}}\right)  =\mathfrak{Be}\left(
P_{j_{1}j_{2}k_{1}k_{2}}\right)  =\sum_{1\leq p_{1}<p_{2},q_{1}<q_{2}\leq
r}\left\vert Q_{p_{1}p_{2}q_{1}q_{2}}\right)  \left(  \left.  Q_{p_{1}%
p_{2}q_{1}q_{2}}\right\vert P_{j_{1}j_{2}k_{1}k_{2}}\right)  ,\nonumber
\end{equation}
which leads to%
\begin{align}
&  Q_{j_{1}j_{2}k_{1}k_{2}}\left(  \mathbf{z}\right)  =\left[  \mathfrak{Be}%
\left(  P_{j_{1}j_{2}k_{1}k_{2}}\right)  \right]  \left(  \mathbf{z}\right)
=\nonumber\\
&  \int\sum_{1\leq p_{1}<p_{2},q_{1}<q_{2}\leq r}\left(  \left\vert
\mathbf{z}\right\rangle \left\langle \mathbf{z}\right\vert \left\vert
Q_{p_{1}p_{2}q_{1}q_{2}}\right.  \right)  \left(  Q_{p_{1}p_{2}q_{1}q_{2}%
}\left\vert \left\vert \mathbf{z}^{\prime}\right\rangle \left\langle
\mathbf{z}^{\prime}\right\vert \right.  \right)  \left(  \left.  \left\vert
\mathbf{z}^{\prime}\right\rangle \left\langle \mathbf{z}^{\prime}\right\vert
\right\vert P_{j_{1}j_{2}k_{1}k_{2}}\right)  d\mu\left(  \mathbf{z}^{\prime
}\right) \nonumber\\
&  =\int\mathfrak{Be}\left(  \mathbf{z,z}^{\prime}\right)  \left(  \left\vert
\mathbf{z}^{\prime}\right\rangle \left\langle \mathbf{z}^{\prime}\right\vert
\left\vert P_{j_{1}j_{2}k_{1}k_{2}}\right.  \right)  d\mu\left(
\mathbf{z}^{\prime}\right)  .
\end{align}
Thus
\begin{equation}
\mathfrak{Be}\left(  \mathbf{z,z}^{\prime}\right)  =\sum_{1\leq p_{1}%
<p_{2},q_{1}<q_{2}\leq r}\left(  \left\vert \mathbf{z}\right\rangle
\left\langle \mathbf{z}\right\vert \left\vert Q_{p_{1}p_{2}q_{1}q_{2}}\right.
\right)  \left(  \left.  Q_{p_{1}p_{2}q_{1}q_{2}}\right\vert \left\vert
\mathbf{z}^{\prime}\right\rangle \left\langle \mathbf{z}^{\prime}\right\vert
\right)  =\left\vert \left\langle \mathbf{z}\left\vert \mathbf{z}^{\prime
}\right.  \right\rangle \right\vert ^{2}.
\end{equation}


The Kernel of the inverse Berezin Transformation that transform P-symbols to
Q-symbols
\begin{equation}
\mathfrak{Be}^{-1}:\mathfrak{F}\rightarrow\mathfrak{F}%
\end{equation}
can be obtained by using the identity
\begin{equation}
\left\vert P_{j_{1}j_{2}k_{1}k_{2}}\right)  =\mathfrak{Be}^{-1}\left(
Q_{j_{1}j_{2}k_{1}k_{2}}\right)  =\sum_{1\leq p_{1}<p_{2},q_{1}<q_{2}\leq
r}\left\vert P_{p_{1}p_{2}q_{1}q_{2}}\right)  \left(  \left.  P_{p_{1}%
p_{2}q_{1}q_{2}}\right\vert Q_{j_{1}j_{2}k_{1}k_{2}}\right)
\end{equation}
to give%
\begin{align}
&  P_{j_{1}j_{2}k_{1}k_{2}}\left(  \mathbf{z}\right)  =\left[  \mathfrak{Be}%
^{-1}\left(  Q_{j_{1}j_{2}k_{1}k_{2}}\right)  \right]  \left(  \mathbf{z}%
\right)  =\nonumber\\
&  \int\sum_{1\leq p_{1}<p_{2},q_{1}<q_{2}\leq r}\left(  \left\vert
\mathbf{z}\right\rangle \left\langle \mathbf{z}\right\vert \left\vert
P_{p_{1}p_{2}q_{1}q_{2}}\right.  \right)  \left(  \left.  P_{p_{1}p_{2}%
q_{1}q_{2}}\right\vert \left\vert \mathbf{z}^{\prime}\right\rangle
\left\langle \mathbf{z}^{\prime}\right\vert \right)  \left(  \left\vert
\mathbf{z}^{\prime}\right\rangle \left\langle \mathbf{z}^{\prime}\right\vert
\left\vert Q_{j_{1}j_{2}k_{1}k_{2}}\right.  \right)  d\mu\left(
\mathbf{z}^{\prime}\right) \nonumber\\
&  =\int\sum_{\substack{1\leq j_{1}<j_{2},k_{1}<k_{2}\leq r\\1\leq p_{1}%
<p_{2},q_{1}<q_{2}\leq r}}\left(  \left\vert \mathbf{z}\right\rangle
\left\langle \mathbf{z}\right\vert \left\vert Q_{j_{1}j_{2}k_{1}k_{2}}\right.
\right)  \Delta_{j_{1}j_{2}k_{1}k_{2}p_{1}p_{2}q_{1}q_{2}}^{-2}\left(  \left.
Q_{p_{1}p_{2}q_{1}q_{2}}\right\vert \left\vert \mathbf{z}^{\prime
}\right\rangle \left\langle \mathbf{z}^{\prime}\right\vert \right) \nonumber\\
&  \times\left(  \left\vert \mathbf{z}^{\prime}\right\rangle \left\langle
\mathbf{z}^{\prime}\right\vert \left\vert Q_{j_{1}j_{2}k_{1}k_{2}}\right.
\right)  d\mu\left(  \mathbf{z}^{\prime}\right) \nonumber\\
&  =\int\mathfrak{Be}^{-1}\left(  \mathbf{z,z}^{\prime}\right)  \left(
\left\vert \mathbf{z}^{\prime}\right\rangle \left\langle \mathbf{z}^{\prime
}\right\vert \left\vert Q_{j_{1}j_{2}k_{1}k_{2}}\right.  \right)  d\mu\left(
\mathbf{z}^{\prime}\right)  ,
\end{align}
where
\begin{equation}
\mathfrak{Be}^{-1}\left(  \mathbf{z,z}^{\prime}\right)  =\sum_{\substack{1\leq
j_{1}<j_{2},k_{1}<k_{2}\leq r\\1\leq p_{1}<p_{2},q_{1}<q_{2}\leq r}}\left(
\left\vert \mathbf{z}\right\rangle \left\langle \mathbf{z}\right\vert
\left\vert Q_{j_{1}j_{2}k_{1}k_{2}}\right.  \right)  \Delta_{j_{1}j_{2}%
k_{1}k_{2}p_{1}p_{2}q_{1}q_{2}}^{-2}\left(  \left.  Q_{p_{1}p_{2}q_{1}q_{2}%
}\right\vert \left\vert \mathbf{z}^{\prime}\right\rangle \left\langle
\mathbf{z}^{\prime}\right\vert \right)  .
\end{equation}


$\lceil$The Berezin Transform gives
\begin{equation}
Q_{p_{1}p_{2}q_{1}q_{2}}\left(  \mathbf{z}\right)  =\int P_{p_{1}p_{2}%
q_{1}q_{2}}\left(  \mathbf{z}^{\prime}\right)  \left\vert \left\langle
\mathbf{z}\left\vert \mathbf{z}^{\prime}\right.  \right\rangle \right\vert
^{2}d\mu\left(  \mathbf{z}^{\prime}\right)  ,
\end{equation}
which leads to (using the biorthogonality relationships)
\begin{align}
&  \sum_{1\leq j_{1}<j_{2},k_{1}<k_{2}\leq r}P_{j_{1}j_{2}k_{1}k_{2}}\left(
\mathbf{z}\right)  \Delta_{j_{1}j_{2}k_{1}k_{2}p_{1}p_{2}q_{1}q_{2}}=\int
P_{p_{1}p_{2}q_{1}q_{2}}\left(  \mathbf{z}^{\prime}\right)  \left\vert
\left\langle \mathbf{z}\left\vert \mathbf{z}^{\prime}\right.  \right\rangle
\right\vert ^{2}d\mu\left(  \mathbf{z}^{\prime}\right) \nonumber\\
&  \qquad\qquad\qquad\qquad\qquad\qquad\qquad\left.  \Downarrow\right.
\nonumber\\
&  P_{l_{1}l_{2}m_{1}m_{2}}\left(  \mathbf{z}\right)  =\int\sum_{1\leq
j_{1}<j_{2},k_{1}<k_{2}\leq r}Q_{j_{1}j_{2}k_{1}k_{2}}\left(  \mathbf{z}%
^{\prime}\right)  \Delta_{j_{1}j_{2}k_{1}k_{2}l_{1}l_{2}m_{1}m_{2}}%
^{-2}\left\vert \left\langle \mathbf{z}\left\vert \mathbf{z}^{\prime}\right.
\right\rangle \right\vert ^{2}d\mu\left(  \mathbf{z}^{\prime}\right)
\end{align}
and
\begin{equation}
Q_{p_{1}p_{2}q_{1}q_{2}}\left(  \mathbf{z}\right)  =\int\sum_{1\leq
j_{1}<j_{2},k_{1}<k_{2}\leq r}Q_{j_{1}j_{2}k_{1}k_{2}}\left(  \mathbf{z}%
^{\prime}\right)  \Delta_{j_{1}j_{2}k_{1}k_{2}p_{1}p_{2}q_{1}q_{2}}%
^{-1}\left\vert \left\langle \mathbf{z}\left\vert \mathbf{z}^{\prime}\right.
\right\rangle \right\vert ^{2}d\mu\left(  \mathbf{z}^{\prime}\right)  .
\end{equation}
$\rfloor$

The Q and P-symols of any operator are
\begin{equation}
Q_{A}\left(  \mathbf{z}\right)  =\sum_{1\leq j_{1}<j_{2}\leq r,1\leq
k_{1}<k_{2}\leq r}A_{j_{1}j_{2}k_{1}k_{2}}Q_{j_{1}j_{2}k_{1}k_{2}}\left(
\mathbf{z}\right)  ,
\end{equation}%
\begin{equation}
P_{A}\left(  \mathbf{z}\right)  =\sum_{1\leq j_{1}<j_{2}\leq r,1\leq
k_{1}<k_{2}\leq r}A_{j_{1}j_{2}k_{1}k_{2}}P_{j_{1}j_{2}k_{1}k_{2}}\left(
\mathbf{z}\right)  ,
\end{equation}
where
\begin{equation}
A_{j_{1}j_{2}k_{1}k_{2}}=\int P_{A}\left(  \mathbf{z}\right)  Q_{j_{1}%
j_{2}k_{1}k_{2}}\left(  \mathbf{z}\right)  d\mu\left(  \mathbf{z}\right)
=\int P_{j_{1}j_{2}k_{1}k_{2}}\left(  \mathbf{z}\right)  Q_{A}\left(
\mathbf{z}\right)  d\mu\left(  \mathbf{z}\right)  .
\end{equation}


\subsection{Dual Spaces of P and Q-Symbols}

In finite dimensions the space $\mathfrak{F}$ is spanned by both $\left\{
Q_{j_{1}j_{2}k_{1}k_{2}}\right\}  $ and $\left\{  P_{j_{1}j_{2}k_{1}k_{2}%
}\right\}  $ and is complete in the norms $\left\Vert {}\right\Vert _{\infty
},\left\Vert {}\right\Vert _{1}$ and $\left\Vert {}\right\Vert _{2}$ where
\begin{subequations}
\begin{align}
\left\Vert Q_{A}\right\Vert _{\infty} &  =\sup\left\{  \left\vert Q\left(
\mathbf{z}\right)  \right\vert \right\}  =\sup\left\{  \left\vert \left\langle
\mathbf{z}\left\vert A\mathbf{z}\right.  \right\rangle \right\vert \right\}
,\nonumber\\
\left\Vert Q_{A}\right\Vert _{1} &  =\int\left\vert Q\left(  \mathbf{z}%
\right)  \right\vert d\mu\left(  \mathbf{z}\right)  =\int\left\vert
\left\langle \mathbf{z}\left\vert A\mathbf{z}\right.  \right\rangle
\right\vert d\mu\left(  \mathbf{z}\right)  ,\nonumber\\
\left\Vert Q_{A}\right\Vert _{2} &  =\left\{  \int\left\vert Q\left(
\mathbf{z}\right)  \right\vert ^{2}d\mu\left(  \mathbf{z}\right)  \right\}
^{\frac{1}{2}}=\left\{  \int\left\vert \left\langle \mathbf{z}\left\vert
A\mathbf{z}\right.  \right\rangle \right\vert ^{2}d\mu\left(  \mathbf{z}%
\right)  \right\}  ^{\frac{1}{2}}.
\end{align}
If $A\geq0$ then
\end{subequations}
\begin{equation}
\left\Vert Q_{A}\right\Vert _{1}=\int\left\langle \mathbf{z}\left\vert
A\mathbf{z}\right.  \right\rangle d\mu\left(  \mathbf{z}\right)
=\operatorname*{Tr}\left\{  A\right\}  .
\end{equation}
$\lceil$%
\begin{equation}
\operatorname*{Tr}\left\{  A^{\dagger}A\right\}  =\int\left\langle
\mathbf{z}\left\vert A^{\dagger}A\mathbf{z}\right.  \right\rangle d\mu\left(
\mathbf{z}\right)  \rfloor
\end{equation}


\subsection{Algebraic Structure}

\paragraph{P-Symbol products}%

\begin{align}
&  \int P_{A}\left(  \mathbf{z}_{1}\right)  \left\vert \mathbf{z}%
_{1}\right\rangle \left\langle \mathbf{z}_{1}\right\vert d\mu\left(
\mathbf{z}_{1}\right)  \int P_{A}\left(  \mathbf{z}_{1}\right)  \left\vert
\mathbf{z}_{2}\right\rangle \left\langle \mathbf{z}_{2}\right\vert d\mu\left(
\mathbf{z}_{2}\right) \\
&  =\int P_{A}\left(  \mathbf{z}_{1}\right)  P_{B}\left(  \mathbf{z}%
_{2}\right)  S\left(  \mathbf{\bar{z}}_{1},\mathbf{z}_{2}\right)  \left\vert
\mathbf{z}_{1}\right\rangle \left\langle \mathbf{z}_{2}\right\vert d\mu\left(
\mathbf{z}_{1}\right)  d\mu\left(  \mathbf{z}_{2}\right) \\
&  =\int P_{AB}\left(  \mathbf{z}_{1}\right)  \left\vert \mathbf{z}%
_{1}\right\rangle \left\langle \mathbf{z}_{1}\right\vert d\mu\left(
\mathbf{z}_{1}\right) \\
P_{AB}\left(  \mathbf{z}_{1}\right)   &  =\int P_{A}\left(  \mathbf{z}%
_{1}\right)  S\left(  \mathbf{\bar{z}}_{1},\mathbf{z}_{2}\right)  P_{B}\left(
\mathbf{z}_{2}\right)
\end{align}


\paragraph{Q-Symbol products}

Pointwise product produces a commutative $w^{\ast}$- algebra
\begin{equation}
Q_{A\cdot B}\left(  \mathbf{\bar{z}},\mathbf{z}\right)  =Q_{A}\left(
\mathbf{\bar{z}},\mathbf{z}\right)  Q_{B}\left(  \mathbf{\bar{z}}%
,\mathbf{z}\right)
\end{equation}
this product is distinct from non commutative products such as
\begin{equation}
Q_{A\circ B}\left(  \mathbf{\bar{z}},\mathbf{z}\right)  =\left\langle
\mathbf{z}\left\vert AB\mathbf{z}\right.  \right\rangle ,
\end{equation}%
\begin{equation}
Q_{AB_{+}}\left(  \mathbf{\bar{z}},\mathbf{z}\right)  =\frac{1}{2}\left\langle
\mathbf{z}\left\vert \left(  AB+BA\right)  \mathbf{z}\right.  \right\rangle
=\frac{1}{2}\left(  Q_{A\circ B}\left(  \mathbf{\bar{z}},\mathbf{z}\right)
+Q_{B\circ A}\left(  \mathbf{\bar{z}},\mathbf{z}\right)  \right)  .
\end{equation}


\subsection{Independent Particle Example}

Let the ground state of a non interacting two particle system be
\begin{equation}
D_{IPS}^{2}=\left\vert \varphi_{1}\varphi_{2}\right\rangle \left\langle
\varphi_{1}\varphi_{2}\right\vert
\end{equation}
then%

\begin{align}
Q_{\gamma_{2\varphi}}\left(  \mathbf{z}\right)   &  =\left\langle
\mathbf{z}\left\vert \left\vert \varphi_{1}\varphi_{2}\right\rangle
\left\langle \varphi_{1}\varphi_{2}\right\vert \right.  \mathbf{z}%
\right\rangle =\left\langle \varphi_{1}\varphi_{2}\left\vert \left\vert
\mathbf{z}\right\rangle \left\langle \mathbf{z}\right\vert \right.
\varphi_{1}\varphi_{2}\right\rangle =\gamma_{2\varphi}\left(  \mathbf{z}%
\right)  _{1212}\left\langle \mathbf{z}\left\vert \mathbf{z}\right.
\right\rangle \equiv\gamma_{2}\left(  \mathbf{z}\right)  _{1212}\left\langle
\mathbf{z}\left\vert \mathbf{z}\right.  \right\rangle \nonumber\\
&  =\left\{  \gamma_{1}\left(  \mathbf{z}\right)  _{11}\gamma_{1}\left(
\mathbf{z}\right)  _{22}-\gamma_{1}\left(  \mathbf{z}\right)  _{12}\gamma
_{1}\left(  \mathbf{z}\right)  _{21}\right\}  \left\langle \mathbf{z}%
\left\vert \mathbf{z}\right.  \right\rangle =\left\vert
\begin{array}
[c]{cc}%
\gamma_{1}\left(  \mathbf{z}\right)  _{11} & \gamma_{1}\left(  \mathbf{z}%
\right)  _{12}\\
\gamma_{1}\left(  \mathbf{z}\right)  _{21} & \gamma_{1}\left(  \mathbf{z}%
\right)  _{22}%
\end{array}
\right\vert \left\langle \mathbf{z}\left\vert \mathbf{z}\right.  \right\rangle
\nonumber\\
&  =\left\vert
\begin{array}
[c]{cc}%
\left[  \left(  \mathbb{I}_{2}+\mathbf{z}^{\dagger}\mathbf{z}\right)
^{-1}\right]  _{11} & \left[  \left(  \mathbb{I}_{2}+\mathbf{z}^{\dagger
}\mathbf{z}\right)  ^{-1}\right]  _{12}\\
\left[  \left(  \mathbb{I}_{2}+\mathbf{z}^{\dagger}\mathbf{z}\right)
^{-1}\right]  _{21} & \left[  \left(  \mathbb{I}_{2}+\mathbf{z}^{\dagger
}\mathbf{z}\right)  ^{-1}\right]  _{22}%
\end{array}
\right\vert \left\langle \mathbf{z}\left\vert \mathbf{z}\right.  \right\rangle
=\det\left(  \mathbb{I}_{2}+\mathbf{z}^{\dagger}\mathbf{z}\right)
^{-1}\left\langle \mathbf{z}\left\vert \mathbf{z}\right.  \right\rangle =1
\end{align}
and%
\begin{align}
\left\langle \varphi_{1}\varphi_{2}\left\vert \left\vert \mathbf{z}%
\right\rangle \left\langle \mathbf{z}\right\vert \right.  \varphi_{1}%
\varphi_{2}\right\rangle  &  =\left\langle \varphi_{1}\varphi_{2}\left\vert
\left\vert c\left(  \mathbf{z}\right)  \varphi_{1}\varphi_{2}\right\rangle
\left\langle \varphi_{1}\varphi_{2}\right\vert \right.  c\left(
\mathbf{z}\right)  ^{\dagger}\varphi_{1}\varphi_{2}\right\rangle \nonumber\\
&  =\left\vert \left\langle \varphi_{1}\varphi_{2}\left\vert c\left(
\mathbf{z}\right)  \varphi_{1}\varphi_{2}\right.  \right\rangle \right\vert
^{2}=1.
\end{align}
Note that%
\begin{equation}
\det\left(  \mathbb{I}_{2}+\mathbf{z}^{\dagger}\mathbf{z}\right)  =\det\left(
%
\begin{array}
[c]{cc}%
1+\left(  \mathbf{z}^{\dagger}\mathbf{z}\right)  _{11} & \left(
\mathbf{z}^{\dagger}\mathbf{z}\right)  _{12}\\
\left(  \mathbf{z}^{\dagger}\mathbf{z}\right)  _{21} & 1+\left(
\mathbf{z}^{\dagger}\mathbf{z}\right)  _{22}%
\end{array}
\right)  ,
\end{equation}%
\begin{equation}
\gamma_{1}\left(  \mathbf{z}\right)  _{hh^{\prime}}=\left[  \left(
\mathbb{I}_{2}+\mathbf{z}^{\dagger}\mathbf{z}\right)  ^{-1}\right]
_{hh^{\prime}}%
\end{equation}
and
\begin{align}
1 &  =\left\langle \mathbf{z}\left\vert \left\vert \varphi_{1}\varphi
_{2}\right\rangle \left\langle \varphi_{1}\varphi_{2}\right\vert \right.
\mathbf{z}\right\rangle =\left\vert \left\langle \left.  \mathbf{z}\right\vert
\varphi_{1}\varphi_{2}\right\rangle \right\vert ^{2}=\left\vert \left\langle
\left.  \varphi_{1}\varphi_{2}\right\vert c\left(  \mathbf{z}\right)
\varphi_{1}\varphi_{2}\right\rangle \right\vert ^{2}\nonumber\\
&  =\left\vert \left\langle \varphi_{1}\varphi_{2}\left\vert \exp\left\{
\sum\limits_{\substack{1\leq h\leq2\\3\leq p\leq r}}z_{ph}a_{p}^{\dagger}%
a_{h}\right\}  \varphi_{1}\varphi_{2}\right.  \right\rangle \right\vert ^{2},
\end{align}
thus
\begin{equation}
Q_{\gamma_{2\varphi}}\left(  \mathbf{z}\right)  =1\quad\forall\mathbf{z\in
}\mathbb{C}^{2\left(  r-2\right)  }.
\end{equation}
In this case
\begin{align}
&  1=\int Q_{\gamma_{2\varphi}}\left(  \mathbf{z}\right)  d\mu\left(
\mathbf{z}\right)  =\operatorname*{Tr}\left\{  \left\vert \varphi_{1}%
\varphi_{2}\right\rangle \left\langle \varphi_{1}\varphi_{2}\right\vert
\right\}  \nonumber\\
&  =\int d\mu\left(  \mathbf{z}\right)  =\int\det\left(  \mathbb{I}%
_{2}+\mathbf{z}^{\dagger}\mathbf{z}\right)  ^{-\left(  r+1\right)  }\left(
-2\pi i\right)  ^{-2\left(  r-2\right)  }d\mathbf{\bar{z}}d\mathbf{z.}%
\end{align}
The energy functional is
\begin{equation}
E_{N}\left(  D^{2}\right)  =\int P_{K_{2N}}\left(  \mathbf{z}\right)
Q_{D^{2}}\left(  \mathbf{z}\right)  d\mu\left(  \mathbf{z}\right)  =\int
Q_{K_{2N}}\left(  \mathbf{z}\right)  P_{D^{2}}\left(  \mathbf{z}\right)
d\mu\left(  \mathbf{z}\right)
\end{equation}
$\lceil$%
\begin{equation}
\operatorname*{Tr}\left\{  K_{2N}D\right\}  =\int P_{K_{2N}}\left(
\mathbf{z}\right)  \operatorname*{Tr}\left\{  \left\vert \mathbf{z}%
\right\rangle \left\langle \mathbf{z}\right\vert D\right\}  d\mu\left(
\mathbf{z}\right)  =\int P_{K_{2N}}\left(  \mathbf{z}\right)  Q_{D^{2}}\left(
\mathbf{z}\right)  d\mu\left(  \mathbf{z}\right)  \rfloor
\end{equation}
and the minimum value is%
\begin{equation}
E_{N}\left(  D_{\ast}^{2}\right)  =K_{\varphi1212}=\int P_{K_{22}}\left(
\mathbf{z}\right)  Q_{D_{\ast}^{2}}\left(  \mathbf{z}\right)  d\mu\left(
\mathbf{z}\right)  ,
\end{equation}
where the Q-symbol of the Reduced hamiltonian is given by
\begin{align}
&  Q_{K_{22}}\left(  \mathbf{\bar{z}},\mathbf{z}\right)  =\frac{1}{2}%
\sum_{1\leq j_{1},j_{2},k_{1},k_{2}\leq r}K_{\varphi j_{1}j_{2}k_{1}k_{2}%
}\frac{{}}{\left\langle \mathbf{z}\left\vert \mathbf{z}\right.  \right\rangle
}\frac{{}}{\left\langle \mathbf{z}\left\vert \mathbf{z}\right.  \right\rangle
}\nonumber\\
&  =\frac{1}{2}\sum_{1\leq j_{1},j_{2},k_{1},k_{2}\leq r}\left\{  \eta\left(
\delta_{j_{1}k_{1}}\delta_{j_{2}k_{2}}-\delta_{j_{1}k_{2}}\delta_{j_{2}k_{1}%
}\right)  +\left(  h_{1jk}\delta_{lm}+h_{1lm}\delta_{jk}\right)  \right\}
\nonumber\\
&  \times K_{\varphi j_{1}j_{2}k_{1}k_{2}}\left\{  \gamma_{1}\left(
\mathbf{z}\right)  _{j_{1}k_{1}}\gamma_{1}\left(  \mathbf{z}\right)
_{j_{2}k_{2}}-\gamma_{1}\left(  \mathbf{z}\right)  _{j_{1}k_{2}}\gamma
_{1}\left(  \mathbf{z}\right)  _{j_{2}k_{1}}\right\}  ,
\end{align}
and its P-symbol $P_{K_{22}}\left(  \mathbf{z}\right)  $ has to be found via
the inverse Berezin transform
\begin{equation}
P_{K_{22}}\left(  \mathbf{z}\right)  =\left[  \mathfrak{Be}^{-1}\left(
Q_{K_{22}}\left(  \mathbf{z}\right)  \right)  \right]  \left(  \mathbf{z}%
\right)  =\int\mathfrak{Be}^{-1}\left(  \mathbf{\bar{z},z}^{\prime}\right)
Q_{K_{22}}\left(  \mathbf{z}^{\prime}\right)  d\mu\left(  \mathbf{z}^{\prime
}\right)  ,
\end{equation}
and%
\begin{align}
Q_{D_{\ast}^{2}}\left(  \mathbf{z}\right)   &  =\left[  D_{\ast2}\right]
_{1212}\gamma_{2}\left(  \mathbf{z}\right)  _{1212}\nonumber\\
&  =\left[  D_{\ast2}\right]  _{1212}\left\{  \gamma_{1}\left(  \mathbf{z}%
\right)  _{11}\gamma_{1}\left(  \mathbf{z}\right)  _{22}-\gamma_{1}\left(
\mathbf{z}\right)  _{12}\gamma_{1}\left(  \mathbf{z}\right)  _{21}\right\}
\nonumber\\
&  =\gamma_{1}\left(  \mathbf{z}\right)  _{11}\gamma_{1}\left(  \mathbf{z}%
\right)  _{22}-\gamma_{1}\left(  \mathbf{z}\right)  _{12}\gamma_{1}\left(
\mathbf{z}\right)  _{21}=\det\left(  \mathbb{I}_{2}+\mathbf{z}^{\dagger
}\mathbf{z}\right)  ^{-1}.
\end{align}
The energy is given as
\begin{align}
E\left(  D_{\ast}^{2}\right)   &  =\inf_{D^{2}\in S_{22}=S_{2}}%
\operatorname*{Tr}\left\{  K_{22}D^{2}\right\}  \nonumber\\
&  \parallel\nonumber\\
E\left(  Q_{D_{\ast}^{2}}\right)   &  =\inf_{Q\in\mathfrak{Q}_{12}}\left\{
\int P_{K_{22}}\left(  \mathbf{z}\right)  d\mu\left(  \mathbf{z}\right)
\right\}  \nonumber\\
&  =K_{\varphi1212},
\end{align}
where%
\begin{equation}
K_{22}=\int P_{K_{22}\varphi}\left(  \mathbf{z}\right)  \left\vert c\left(
\mathbf{z}\right)  \varphi_{1}\varphi_{2}\right\rangle \left\langle c\left(
\mathbf{z}\right)  \varphi_{1}\varphi_{2}\right\vert d\mu\left(
\mathbf{z}\right)  .
\end{equation}


\subsection{BiOrthogonal Basis}%

\begin{align}
A\left\vert \mathbf{\varphi}\right\rangle  &  =\left\vert \mathbf{\varphi
}\right\rangle \mathbb{A}\\
\left\langle \left.  \mathbf{\psi}\right\vert \mathbf{\varphi}\right\rangle
&  =\mathbb{I}\\
\Delta_{\psi}  &  =\left\langle \left.  \mathbf{\psi}\right\vert \mathbf{\psi
}\right\rangle \Rightarrow\left\langle \left.  \mathbf{\psi}\right\vert
\mathbf{\psi}\right\rangle \mathbf{\Delta}_{\psi}^{-1}=\left\langle \left.
\mathbf{\psi}\right\vert \Delta_{\psi}^{-1}\mathbf{\psi}\right\rangle
=\mathbb{I}\Rightarrow\left\vert \mathbf{\varphi}\right\rangle =\left\vert
\Delta_{\psi}^{-1}\mathbf{\psi}\right\rangle \\
P  &  =\left\vert \mathbf{\varphi}\right\rangle \mathbf{\Delta}_{\varphi}%
^{-1}\left\langle \mathbf{\varphi}\right\vert =\left\vert \mathbf{\psi
}\right\rangle \mathbf{\Delta}_{\psi}^{-1}\left\langle \mathbf{\psi
}\right\vert =\left\vert \mathbf{\varphi}\right\rangle \left\langle
\mathbf{\psi}\right\vert =\left\vert \mathbf{\varphi}\right\rangle
\left\langle \mathbf{\psi}\right\vert
\end{align}


\section{Extra}%

\begin{equation}
\sum_{1\leq p_{1}<p_{2},q_{1}<q_{2}\leq r}P_{p_{1}p_{2}q_{1}q_{2}}\left(
\mathbf{z}^{\prime}\right)  \Delta_{p_{1}p_{2}q_{1}q_{2}j_{1}j_{2}k_{1}k_{2}%
}^{-1}%
\end{equation}
%

\begin{align}
&  =\int\Upsilon_{j_{1}j_{2}k_{1}k_{2}}\left(  \mathbf{z}^{\prime}\right)
\left\vert \left\langle \mathbf{z}\left\vert \mathbf{z}^{\prime}\right.
\right\rangle \right\vert ^{2}d\mu\left(  \mathbf{z}^{\prime}\right)
\nonumber\\
\Upsilon_{j_{1}j_{2}k_{1}k_{2}}\left(  \mathbf{z}^{\prime}\right)   &
=\sum_{1\leq p_{1}<p_{2},q_{1}<q_{2}\leq r}P_{p_{1}p_{2}q_{1}q_{2}}\left(
\mathbf{z}^{\prime}\right)  \Delta_{p_{1}p_{2}q_{1}q_{2}j_{1}j_{2}k_{1}k_{2}%
}^{-1}.
\end{align}
this gives%
\begin{equation}
\Upsilon_{j_{1}j_{2}k_{1}k_{2}}\left(  \mathbf{z}^{\prime}\right)
=\sum_{1\leq p_{1}<p_{2},q_{1}<q_{2}\leq r}Q_{p_{1}p_{2}q_{1}q_{2}}\left(
\mathbf{z}\right)  \Delta_{p_{1}p_{2}q_{1}q_{2}j_{1}j_{2}k_{1}k_{2}}^{-2}%
\end{equation}
and
\begin{equation}
P_{j_{1}j_{2}k_{1}k_{2}}\left(  \mathbf{z}\right)  =\int\sum_{1\leq
p_{1}<p_{2},q_{1}<q_{2}\leq r}Q_{p_{1}p_{2}q_{1}q_{2}}\left(  \mathbf{z}%
\right)  \Delta_{p_{1}p_{2}q_{1}q_{2}j_{1}j_{2}k_{1}k_{2}}^{-2}\left\vert
\left\langle \mathbf{z}\left\vert \mathbf{z}^{\prime}\right.  \right\rangle
\right\vert ^{2}d\mu\left(  \mathbf{z}^{\prime}\right)  .
\end{equation}
As%
\begin{equation}
Q_{A}\left(  \mathbf{z}\right)  =\int P_{A}\left(  \mathbf{z}^{\prime}\right)
\left\vert \left\langle \mathbf{z}\left\vert \mathbf{z}^{\prime}\right.
\right\rangle \right\vert ^{2}d\mu\left(  \mathbf{z}^{\prime}\right)
\end{equation}
one obtains%
\begin{align}
P_{A}\left(  \mathbf{z}\right)   &  =\int\sum_{1\leq p_{1}<p_{2},q_{1}%
<q_{2}\leq r}A_{j_{1}j_{2}k_{1}k_{2}}P_{p_{1}p_{2}q_{1}q_{2}}\left(
\mathbf{z}^{\prime}\right)  \Delta_{p_{1}p_{2}q_{1}q_{2}j_{1}j_{2}k_{1}k_{2}%
}^{-1}\left\vert \left\langle \mathbf{z}\left\vert \mathbf{z}^{\prime}\right.
\right\rangle \right\vert ^{2}d\mu\left(  \mathbf{z}^{\prime}\right)
\nonumber\\
&  =\int\Upsilon_{A}\left(  \mathbf{z}^{\prime}\right)  \left\vert
\left\langle \mathbf{z}\left\vert \mathbf{z}^{\prime}\right.  \right\rangle
\right\vert ^{2}d\mu\left(  \mathbf{z}^{\prime}\right) \nonumber\\
\Upsilon_{A}\left(  \mathbf{z}^{\prime}\right)   &  =\sum_{\substack{1\leq
p_{1}<p_{2},q_{1}<q_{2}\leq r\\1\leq j_{1}<j_{2},k_{1}<k_{2}\leq r}%
}A_{j_{1}j_{2}k_{1}k_{2}}P_{p_{1}p_{2}q_{1}q_{2}}\left(  \mathbf{z}^{\prime
}\right)  \Delta_{p_{1}p_{2}q_{1}q_{2}j_{1}j_{2}k_{1}k_{2}}^{-1}.\rfloor
\end{align}
As%
\begin{equation}
P_{p_{1}p_{2}q_{1}q_{2}}\left(  \mathbf{z}\right)  =\sum_{1\leq j_{1}%
<j_{2},k_{1}<k_{2}\leq r}Q_{j_{1}j_{2}k_{1}k_{2}}\left(  \mathbf{z}\right)
\Delta_{j_{1}j_{2}k_{1}k_{2}p_{1}p_{2}q_{1}q_{2}}^{-1}\nonumber
\end{equation}


\subsection{}


\end{document}